\documentclass[11pt]{article}
\usepackage{fontspec}
\setmainfont[
  Path=fonts/,
  Extension=.ttf,
  UprightFont=*-400,
  BoldFont=*-700
]{NotoSerifSC}
\setsansfont[
  Path=fonts/,
  Extension=.ttf,
  UprightFont=*-400,
  BoldFont=*-700
]{NotoSerifSC}
\XeTeXlinebreaklocale "zh"
\XeTeXlinebreakskip=0pt plus 1pt
\usepackage[a4paper,margin=1.70cm]{geometry}
\usepackage{amsmath,amssymb,mathtools,bm}
\usepackage{booktabs,longtable,array}
\usepackage{enumitem}
\setlength{\parindent}{2em}
\setlength{\parskip}{0.35em}
\setlist{nosep,leftmargin=2em}
\allowdisplaybreaks[3]
\numberwithin{equation}{section}
\pagestyle{plain}

\newcommand{\R}{\mathbb{R}}
\newcommand{\E}{\mathbb{E}}
\newcommand{\Pp}{\mathbb{P}}
\newcommand{\1}{\mathbf{1}}
\newcommand{\T}{\mathsf{T}}
\newcommand{\F}{\mathrm{F}}
\newcommand{\cN}{\mathcal{N}}
\newcommand{\cL}{\mathcal{L}}
\newcommand{\cS}{\mathcal{S}}
\newcommand{\diag}{\operatorname{diag}}
\newcommand{\tr}{\operatorname{tr}}
\newcommand{\Var}{\operatorname{Var}}
\newcommand{\Cov}{\operatorname{Cov}}
\newcommand{\TopK}{\operatorname{TopK}}
\newcommand{\sgn}{\operatorname{sign}}

\begin{document}
\begin{center}
  {\LARGE\bfseries 29\_FlyLoRA}
\end{center}
\vspace{0.8em}

\section{符号与维度}

\begin{longtable}{>{$}l<{$} >{$}l<{$} p{10cm}}
\toprule
\text{符号} & \text{维度} & \text{含义}\\
\midrule
W_0 & \R^{m\times n} & 冻结的预训练权重。\\
x & \R^n & 单个输入 token 的列向量。\\
A & \R^{r\times n} & 下投影；FlyLoRA 中稀疏、随机初始化并冻结。\\
B & \R^{m\times r} & 上投影；FlyLoRA 中唯一通过梯度训练的适配器矩阵。\\
a_i=A[i,:] & \R^{1\times n} & $A$ 的第 $i$ 行。\\
b_i=B[:,i] & \R^{m\times 1} & $B$ 的第 $i$ 列。\\
y=Ax & \R^r & 稀疏投影后的 rank-wise 激活。\\
\Lambda & \R^{r\times r} & top-$k$ 二值对角掩码。\\
\alpha/r & \R & LoRA 缩放系数，以下简记为 $c=\alpha/r$。\\
p & \{1,\ldots,n-1\} & 每行非零位置数；稀疏率 $q=\rho=p/n$。\\
k & \{1,\ldots,r\} & 每个样本激活的 rank-1 专家数。\\
\delta & \R^m & 输出端反向传播向量 $\partial\cL/\partial h$。\\
\bottomrule
\end{longtable}

所有向量默认是列向量；$\|\cdot\|_2$ 对向量表示欧氏范数、对矩阵表示谱范数，$\|\cdot\|_\F$ 表示 Frobenius 范数。

\section{从 LoRA 到 rank-wise 专家}

\subsection{LoRA 的基本分解}

LoRA 用秩不超过 $r$ 的矩阵模拟全量更新：
\begin{equation}
  W'=W_0+\Delta W,
  \qquad
  \Delta W=cBA,
  \qquad c=\frac{\alpha}{r}.
  \label{eq:lora-weight}
\end{equation}
维度逐项检查为
\begin{equation}
  (m\times r)(r\times n)=m\times n,
  \qquad
  (m\times n)(n\times 1)=m\times 1.
\end{equation}
因此前向传播是
\begin{equation}
  f_{\mathrm{LoRA}}(x)
  =W'x
  =W_0x+cBAx.
  \label{eq:lora-forward}
\end{equation}
由于
\begin{equation}
  B=[b_1,\ldots,b_r],
  \qquad
  A=\begin{bmatrix}a_1\\ \vdots\\ a_r\end{bmatrix},
\end{equation}
矩阵乘积可按中间指标展开。对任意 $\mu\in\{1,\ldots,m\}$、$\nu\in\{1,\ldots,n\}$，
\begin{align}
  (BA)_{\mu\nu}
  &=\sum_{i=1}^{r}B_{\mu i}A_{i\nu} \\
  &=\sum_{i=1}^{r}(b_i)_\mu(a_i)_\nu
  =\left(\sum_{i=1}^{r}b_i a_i\right)_{\mu\nu}.
\end{align}
故有精确恒等式
\begin{equation}
  BA=\sum_{i=1}^{r}b_i a_i,
  \qquad
  BAx=\sum_{i=1}^{r}b_i(a_i x).
  \label{eq:rank-one-sum}
\end{equation}
其中 $b_i a_i\in\R^{m\times n}$ 的秩至多为 $1$，而 $a_i x\in\R$ 是标量。于是每个中间维度都可以看成一个 rank-1 专家：
\begin{equation}
  E_i(x)=b_i(a_i x)\in\R^m.
\end{equation}

\subsection{SVD 类比的准确含义}

若对某个固定更新矩阵做奇异值分解，
\begin{equation}
  \Delta W=U\Sigma V^{\mathsf{T}}
  =\sum_{i=1}^{s}\sigma_i u_i v_i^{\mathsf{T}},
  \qquad
  \sigma_1\ge\cdots\ge\sigma_s\ge 0,
\end{equation}
则 Eckart--Young--Mirsky 定理给出
\begin{equation}
  \underset{\operatorname{rank}(M)\le k}{\min}
  \|\Delta W-M\|_\F^2
  =\sum_{i=k+1}^{s}\sigma_i^2,
\end{equation}
最优解为 $\sum_{i=1}^{k}\sigma_i u_i v_i^{\mathsf{T}}$。FlyLoRA 借用的是“保留较重要 rank-1 分量”的直觉，但必须注意：一般的 $b_i a_i$ 并不是 $\Delta W$ 的奇异分量，$|a_i x|$ 也不是奇异值。因此，Eckart--Young--Mirsky 定理\textbf{不能直接证明}按 $|a_i x|$ 选 top-$k$ 就是 $BAx$ 的最优近似。

对单个专家，贡献范数有精确分解
\begin{equation}
  \|b_i a_i x\|_2
  =\|b_i(a_i x)\|_2
  =|a_i x|\,\|b_i\|_2.
  \label{eq:expert-magnitude}
\end{equation}
因此按 $|a_i x|$ 排序等同于按真实贡献范数排序，当且仅当各列 $\|b_i\|_2$ 相同；若
\begin{equation}
  0<\beta_{\min}\le \|b_i\|_2\le\beta_{\max},
\end{equation}
则有
\begin{equation}
  \beta_{\min}|a_i x|
  \le \|b_i a_i x\|_2
  \le \beta_{\max}|a_i x|.
\end{equation}
这说明代理分数的失真至多受列范数比 $\beta_{\max}/\beta_{\min}$ 控制。

令 $z=Ax$，$S$ 是 $|z_i|$ 最大的 $k$ 个坐标，截断误差为
\begin{equation}
  e_S=B(I-\Lambda_S)z=\sum_{i\notin S}b_i z_i.
\end{equation}
总能得到
\begin{align}
  \|e_S\|_2
  &\le\sum_{i\notin S}\|b_i\|_2|z_i| \\
  &\le\beta_{\max}\sum_{i\notin S}|z_i| \\
  &\le\beta_{\max}\sqrt{r-k}\,\|z_{S^c}\|_2.
  \label{eq:truncation-bound}
\end{align}
若进一步有 $B^{\mathsf{T}} B=\beta^2I_r$，则
\begin{equation}
  \|e_S\|_2^2
  =z^{\mathsf{T}}(I-\Lambda_S)B^{\mathsf{T}} B(I-\Lambda_S)z
  =\beta^2\sum_{i\notin S}|z_i|^2,
\end{equation}
此时保留最大的 $|z_i|$ 才是所有坐标子集中的严格最优解。

\section{FlyLoRA 前向传播、路由和反向传播}

\subsection{稀疏投影与两种 top-$k$ 解释}

FlyLoRA 固定 $A$，每行恰好保留 $p$ 个位置，并在这些位置采样
\begin{equation}
  A_{ij}\mid(j\in S_i)\sim\cN\!\left(0,\frac{1}{r^2}\right),
  \qquad |S_i|=p.
  \label{eq:exact-p-model}
\end{equation}
令 $y=Ax$。论文叙述和式（7）强调按幅值选择，即
\begin{equation}
  \cS_k^{\mathrm{mag}}(y)
  =\{\text{$|y_i|$ 最大的 $k$ 个指标}\}.
  \label{eq:magnitude-topk}
\end{equation}
但原文式（10）加入负载均衡偏置后写成对 $(Ax+d)_i$ 本身逐次取最大值，没有绝对值。对应集合为
\begin{equation}
  \cS_k^{\mathrm{raw}}(x,d)
  =\{\text{$(Ax+d)_i$ 最大的 $k$ 个指标}\}.
  \label{eq:raw-topk}
\end{equation}
两者并不等价。下文凡涉及“贡献幅值”均采用式 \eqref{eq:magnitude-topk}；只涉及二值掩码和梯度稀疏性时，两种规则都适用。

对选中集合 $S$ 定义
\begin{equation}
  \lambda_i=\1\{i\in S\},
  \qquad
  \Lambda=\diag(\lambda_1,\ldots,\lambda_r),
  \qquad
  y'=\Lambda y.
\end{equation}
于是严格且维度一致的前向式为
\begin{align}
  h=f_{\mathrm{FlyLoRA}}(x)
  &=W_0x+cB\Lambda Ax \\
  &=W_0x+c\sum_{i=1}^{r}\lambda_i b_i(a_i x).
  \label{eq:fly-forward}
\end{align}
原文式（8）把分段条件写在 $[By']_i$ 上，但 $By'\in\R^m$，其指标属于输出维度而不是专家维度；维度一致的分段定义应写成
\begin{equation}
  y_i'=\begin{cases}
    y_i,&i\in S,\\
    0,&i\notin S.
  \end{cases}
\end{equation}

\subsection{无损负载均衡偏置}

原文对偏置 $d\in\R^r$ 做非梯度更新：
\begin{equation}
  d_i\leftarrow d_i+u\,\sgn(\bar c_i-c_i).
  \label{eq:bias-update}
\end{equation}
若专家 $i$ 的实际计数 $c_i$ 小于目标计数 $\bar c_i$，则 $d_i$ 增加，使其以后更容易进入 top-$k$；反之则降低。$d$ 不进入专家输出的数值，只改变离散选择集合，所以这种机制不向任务损失增加额外可训练参数。

\subsection{单样本梯度的逐元素推导}

固定本次前向得到的掩码 $\Lambda$，定义
\begin{equation}
  u=\Lambda Ax\in\R^r,
  \qquad
  h=W_0x+cBu,
  \qquad
  \delta=\frac{\partial\cL}{\partial h}\in\R^m.
\end{equation}
对 $B_{\mu i}$，
\begin{align}
  h_\nu
  &=(W_0x)_\nu+c\sum_{j=1}^{r}B_{\nu j}u_j,\\
  \frac{\partial h_\nu}{\partial B_{\mu i}}
  &=c\,\1\{\nu=\mu\}u_i,\\
  \frac{\partial\cL}{\partial B_{\mu i}}
  &=\sum_{\nu=1}^{m}
    \frac{\partial\cL}{\partial h_\nu}
    \frac{\partial h_\nu}{\partial B_{\mu i}}
  =c\,\delta_\mu u_i.
\end{align}
把所有元素重新组成矩阵，得到
\begin{equation}
  \boxed{
  \frac{\partial\cL}{\partial B}
  =c\,\delta u^{\mathsf{T}}
  =c\,\delta (Ax)^{\mathsf{T}}\Lambda.}
  \label{eq:B-gradient}
\end{equation}
第 $i$ 列梯度为
\begin{equation}
  \boxed{
  g_i:=\frac{\partial\cL}{\partial b_i}
  =c\lambda_i(a_i x)\delta.}
  \label{eq:column-gradient}
\end{equation}
若定义不加 top-$k$ 时的稠密梯度
\begin{equation}
  g_i^{(0)}=c(a_i x)\delta,
\end{equation}
则 $g_i=\lambda_i g_i^{(0)}$。这就是附录式（14）
\begin{equation}
  \frac{\partial\cL}{\partial\widetilde B}
  =\frac{\partial\cL}{\partial B}\Lambda
\end{equation}
的逐列含义。

对批量 $\{x_s\}_{s=1}^{N}$，令相应量为 $\delta_s,\Lambda_s$，则
\begin{equation}
  \frac{\partial\cL_{\mathrm{batch}}}{\partial B}
  =c\sum_{s=1}^{N}\delta_s(Ax_s)^{\mathsf{T}}\Lambda_s.
\end{equation}
因为 $A$ 冻结，算法不对离散 top-$k$ 求导；$d$ 由式 \eqref{eq:bias-update} 更新。

\section{稀疏随机投影}

\subsection{两个不同的随机矩阵模型}

主文采用“每行恰好 $p$ 个非零元”的模型 \eqref{eq:exact-p-model}。附录 A.1 为了便于分析，实际改用逐元素独立的替代模型
\begin{equation}
  A_{ij}=\frac{\xi_{ij}G_{ij}}{r},
  \qquad
  \xi_{ij}\sim\operatorname{Bernoulli}(q),
  \quad
  G_{ij}\sim\cN(0,1),
  \quad q=\frac{p}{n},
  \label{eq:bernoulli-gaussian}
\end{equation}
并假设所有 $\xi_{ij},G_{ij}$ 相互独立。两种模型的单个元素边缘分布相同，但前者同一行的支持位置互相依赖；所以从替代模型得到的定理不能不加说明地等同于原始模型定理。

令
\begin{equation}
  z=x-y,
  \qquad
  \sigma^2=\E[A_{ij}^2]=\frac{q}{r^2}=\frac{p}{nr^2}.
  \label{eq:sigma}
\end{equation}
因为各行独立且零均值，
\begin{align}
  \E\|Az\|_2^2
  &=\sum_{u=1}^{r}\E\left(\sum_{j=1}^{n}A_{uj}z_j\right)^2\\
  &=\sum_{u=1}^{r}\sum_{j=1}^{n}\E[A_{uj}^2]z_j^2\\
  &=r\sigma^2\|z\|_2^2.
\end{align}
所以无偏归一化必须是
\begin{equation}
  D_A(z)=\frac{1}{r\sigma^2}\|Az\|_2^2
  =\frac{nr}{p}\|Az\|_2^2,
  \qquad
  \E D_A(z)=\|z\|_2^2.
  \label{eq:normalized-distance}
\end{equation}

\subsection{原论文声称的概率界}

论文定理 3.1/A.2 声称：对任意 $\epsilon>0$，
\begin{align}
 &\Pp\left(
 (1-\epsilon)\|z\|_2^2
 \le D_A(z)
 \le(1+\epsilon)\|z\|_2^2
 \right)\nonumber\\
 &\qquad\ge
 1-\exp\!\left[-(\epsilon^2-\epsilon^3)\frac r4\right]
 -\exp\!\left[
 -\frac{(\epsilon^2-\epsilon^3)r}{2(3p/n+1)}
 \right].
 \label{eq:paper-jl-claim}
\end{align}
其意图是分别控制上尾和下尾，再用
\begin{equation}
  \Pp(L\cap U)
  =1-\Pp(L^c\cup U^c)
  \ge 1-\Pp(L^c)-\Pp(U^c)
\end{equation}
合并两项。

\subsection{附录所引用矩定理的两个尾界}

为了看清常数如何进入结论，把论文附录引用的定理 A.1 改写到当前符号。设随机矩阵 $H\in\R^{r\times n}$ 的元素独立、关于零对称，且
\begin{equation}
  \E[H_{ij}^2]=\sigma^2>0,
  \qquad
  C=\E[H_{ij}^4]<\infty.
\end{equation}
对固定的非零 $z\in\R^n$，其下尾界写为
\begin{equation}
  \Pp\left(
  \frac1r\|Hz\|_2^2
  \le\sigma^2(1-\epsilon)\|z\|_2^2
  \right)
  \le
  \exp\!\left[
  -\frac{(\epsilon^2-\epsilon^3)r}
  {2(C/\sigma^4+1)}
  \right].
  \label{eq:quoted-lower-tail}
\end{equation}
若还存在与阶数无关的 $L>0$，使得对所有整数 $s\ge1$，
\begin{equation}
  \E[H_{ij}^{2s}]
  \le\sigma^{2s}\frac{(2s)!}{2^ss!}L^{2s},
  \label{eq:quoted-moment-condition}
\end{equation}
则引用定理给出的上尾界是
\begin{equation}
  \Pp\left(
  \frac1r\|Hz\|_2^2
  \ge\sigma^2(1+\epsilon)L^2\|z\|_2^2
  \right)
  \le
  \exp\!\left[-(\epsilon^2-\epsilon^3)\frac r4\right].
  \label{eq:quoted-upper-tail}
\end{equation}
注意上尾阈值含有 $L^2$。若 $L=1$，上下尾都围绕期望尺度 $\sigma^2\|z\|_2^2$；若 $L>1$，式 \eqref{eq:quoted-upper-tail} 只控制更远处的尾部。后面将看到，稀疏 Bernoulli--Gaussian 元素需要 $L\ge q^{-1/2}$，这正是原证明无法直接得到对称 $1\pm\epsilon$ 界的关键。

\subsection{偶数阶矩的完整计算}

对任意整数 $s\ge1$，由 $\xi^{2s}=\xi$ 和独立性，
\begin{align}
  \E[A_{ij}^{2s}]
  &=\E\left[\xi_{ij}^{2s}\right]
    \E\left[\left(\frac{G_{ij}}r\right)^{2s}\right]\\
  &=q\,\frac{\E[G^{2s}]}{r^{2s}}.
  \label{eq:moment-step1}
\end{align}
标准高斯偶数阶矩为
\begin{align}
  \E[G^{2s}]
  &=(2s-1)!!\\
  &=(2s-1)(2s-3)\cdots3\cdot1\\
  &=\frac{(2s)!}{2^s s!}.
\end{align}
因此
\begin{equation}
  \boxed{
  \E[A_{ij}^{2s}]
  =q\frac{(2s)!}{2^s s!}\frac1{r^{2s}}.}
  \label{eq:correct-even-moment}
\end{equation}
特别地，
\begin{equation}
  C:=\E[A_{ij}^4]
  =\frac{3q}{r^4}
  =\frac{3p}{nr^4}.
  \label{eq:fourth-moment}
\end{equation}

\subsection{原证明不能直接闭合的四个位置}

\begin{enumerate}[label=\arabic*.]
  \item \textbf{随机模型发生了替换。}“每行恰好 $p$ 个”不等于逐元素独立 Bernoulli$(p/n)$；后者每行非零数是随机变量。
  \item \textbf{高斯矩的幂次。}若方差为 $1/r^2$，则
  \begin{equation}
    \E[(G/r)^{2s}]=(2s-1)!!\,(1/r^2)^s,
  \end{equation}
  而不是 $(2s-1)!!\,(1/r^2)^{2s}$。
  \item \textbf{四阶矩比值。}由 \eqref{eq:sigma}、\eqref{eq:fourth-moment}，
  \begin{equation}
    \frac{C}{\sigma^4}
    =\frac{3q/r^4}{q^2/r^4}
    =\frac3q
    =\frac{3n}{p},
    \label{eq:C-ratio}
  \end{equation}
  不是 $3p/n$。因此直接代入附录定理 A.1 的下尾常数，会得到 $2(3n/p+1)$，而不是式 \eqref{eq:paper-jl-claim} 中的 $2(3p/n+1)$。
  \item \textbf{统一矩常数 $L$。}附录引用的上尾条件要求
  \begin{equation}
    \E[A_{ij}^{2s}]
    \le \sigma^{2s}\frac{(2s)!}{2^s s!}L^{2s}
    \qquad(\forall s\ge1).
  \end{equation}
  代入 \eqref{eq:correct-even-moment} 和 $\sigma^{2s}=q^s/r^{2s}$，等价于
  \begin{equation}
    q\le q^sL^{2s}
    \quad\Longleftrightarrow\quad
    L\ge q^{-\frac{s-1}{2s}}.
  \end{equation}
  对所有 $s$ 统一成立的简单选择是
  \begin{equation}
    L=q^{-1/2}=\sqrt{\frac np}.
    \label{eq:valid-L}
  \end{equation}
  但该引用定理的上尾阈值会包含 $L^2=n/p$，因而不能推出围绕 $1+\epsilon$ 的对称界。
\end{enumerate}

所以，式 \eqref{eq:paper-jl-claim} 应理解为论文的目标性结论；按照附录当前写出的代数，尚不能把它视作已经严格证明的 Johnson--Lindenstrauss 型指数界。

\subsection{一条可完全验证的修正界}

下面仍在逐元素独立模型 \eqref{eq:bernoulli-gaussian} 下，推导一条常数透明、无需引用外部定理的界。令第 $u$ 行投影为
\begin{equation}
  S_u=\sum_{j=1}^{n}A_{uj}z_j,
  \qquad
  \|Az\|_2^2=\sum_{u=1}^{r}S_u^2.
\end{equation}
对独立、零均值、同方差的随机变量展开四次方，只保留“同一指标出现四次”和“两个指标各出现两次”的项：
\begin{align}
  \E[S_u^4]
  &=\sum_{j=1}^{n}Cz_j^4
    +6\sum_{j<\ell}\sigma^4z_j^2z_\ell^2\\
  &=3\sigma^4\left(\sum_{j=1}^{n}z_j^2\right)^2
    +(C-3\sigma^4)\sum_{j=1}^{n}z_j^4.
  \label{eq:row-fourth}
\end{align}
又因为 $\E[S_u^2]=\sigma^2\|z\|_2^2$，
\begin{align}
  \Var(S_u^2)
  &=\E[S_u^4]-(\E[S_u^2])^2\\
  &=2\sigma^4\|z\|_2^4
   +(C-3\sigma^4)\|z\|_4^4.
\end{align}
各行独立，故
\begin{equation}
  \Var(\|Az\|_2^2)
  =r\left[
    2\sigma^4\|z\|_2^4
    +(C-3\sigma^4)\|z\|_4^4
  \right].
\end{equation}
定义集中度参数
\begin{equation}
  \kappa(z)=\frac{\|z\|_4^4}{\|z\|_2^4},
  \qquad
  \frac1n\le\kappa(z)\le1.
\end{equation}
对归一化距离 $R_A(z)=D_A(z)/\|z\|_2^2$，由 \eqref{eq:C-ratio} 得
\begin{align}
  \Var(R_A(z))
  &=\frac{2}{r}
    +\frac1r\left(\frac C{\sigma^4}-3\right)\kappa(z)\\
  &=\frac1r\left[2+\frac{3(1-q)}q\kappa(z)\right].
  \label{eq:normalized-variance}
\end{align}
应用切比雪夫不等式，得到严格有效的双侧界
\begin{equation}
  \boxed{
  \Pp\left(
  \left|\frac{D_A(z)}{\|z\|_2^2}-1\right|\ge\epsilon
  \right)
  \le
  \frac{2+3(1-q)\kappa(z)/q}{r\epsilon^2}.}
  \label{eq:valid-chebyshev}
\end{equation}
式 \eqref{eq:valid-chebyshev} 是可直接验证的多项式尾界，不是论文声称的指数尾界。

\section{top-$k$ 如何降低列间梯度协方差}

\subsection{均匀无放回激活的精确概率}

附录假设每个样本从 $r$ 列中均匀、无放回地选出恰好 $k$ 列。于是对任意 $i\ne j$，
\begin{equation}
  q_1:=\Pp(\lambda_i=1)
  =\frac{\binom{r-1}{k-1}}{\binom rk}
  =\frac kr,
  \label{eq:q1}
\end{equation}
而共同激活概率为
\begin{align}
  q_2&:=\Pp(\lambda_i=1,\lambda_j=1)\\
  &=\frac{\binom{r-2}{k-2}}{\binom rk}
  =\frac{k(k-1)}{r(r-1)}.
  \label{eq:q2}
\end{align}
近似式
\begin{equation}
  q_2\approx\frac{k^2}{r^2}
\end{equation}
只在忽略减一项时成立；精确比值为
\begin{equation}
  \frac{q_2}{k^2/r^2}
  =\frac{r(k-1)}{k(r-1)}.
\end{equation}
当 $k=1$ 时，$q_2=0$，而 $k^2/r^2=1/r^2$，所以“几乎为零”可加强为“在该随机掩码模型下严格为零”。

\subsection{零均值与独立假设下的精确证明}

令 $g_i^{(0)},g_j^{(0)}\in\R^m$ 为不加掩码的随机梯度列，令
\begin{equation}
  \widetilde g_i=\lambda_i g_i^{(0)},
  \qquad
  \widetilde g_j=\lambda_j g_j^{(0)}.
\end{equation}
使用两个关键假设：
\begin{equation}
  (\lambda_i,\lambda_j)\ \text{与}\ (g_i^{(0)},g_j^{(0)})\ \text{独立},
  \qquad
  \E[g_i^{(0)}]=\E[g_j^{(0)}]=0.
  \label{eq:cov-assumptions}
\end{equation}
定义标量型列间协方差
\begin{equation}
  \Sigma_{ij}
  =\E[(g_i^{(0)})^{\mathsf{T}}g_j^{(0)}]
  -\E[g_i^{(0)}]^{\mathsf{T}}\E[g_j^{(0)}].
\end{equation}
在零均值下，$\Sigma_{ij}=\E[(g_i^{(0)})^{\mathsf{T}}g_j^{(0)}]$。于是
\begin{align}
  \widetilde\Sigma_{ij}
  &=\E[\widetilde g_i^{\mathsf{T}}\widetilde g_j]\\
  &=\E[\lambda_i\lambda_j
        (g_i^{(0)})^{\mathsf{T}}g_j^{(0)}]\\
  &=\E[\lambda_i\lambda_j]
    \E[(g_i^{(0)})^{\mathsf{T}}g_j^{(0)}]\\
  &=q_2\Sigma_{ij}.
\end{align}
故精确结论是
\begin{equation}
  \boxed{
  \widetilde\Sigma_{ij}
  =\frac{k(k-1)}{r(r-1)}\Sigma_{ij},
  \qquad i\ne j.}
  \label{eq:exact-cov-reduction}
\end{equation}
论文定理 3.3 的 $k^2/r^2$ 是式 \eqref{eq:exact-cov-reduction} 的渐近近似。

对角项不同：
\begin{align}
  \widetilde\Sigma_{ii}
  &=\E[\lambda_i^2\|g_i^{(0)}\|_2^2]\\
  &=\E[\lambda_i]\E[\|g_i^{(0)}\|_2^2]\\
  &=\frac kr\Sigma_{ii}.
  \label{eq:diag-cov}
\end{align}
所以非对角协方差按约 $(k/r)^2$ 缩小，而方差只按 $k/r$ 缩小。若进一步用相关系数归一化，
\begin{equation}
  \widetilde{\operatorname{Corr}}_{ij}
  =\frac{q_2}{q_1}
   \operatorname{Corr}_{ij}
  =\frac{k-1}{r-1}\operatorname{Corr}_{ij}.
  \label{eq:corr-reduction}
\end{equation}
这比只看协方差更直接地解释论文热力图中非对角相关性的变浅。

\subsection{非零均值时多出的项}

令 $\mu_i=\E[g_i^{(0)}]$、$\mu_j=\E[g_j^{(0)}]$，仍假设掩码与梯度独立。则
\begin{align}
  \widetilde\Sigma_{ij}
  &=q_2\E[(g_i^{(0)})^{\mathsf{T}}g_j^{(0)}]
    -q_1^2\mu_i^{\mathsf{T}}\mu_j\\
  &=q_2\Sigma_{ij}
    +(q_2-q_1^2)\mu_i^{\mathsf{T}}\mu_j.
  \label{eq:nonzero-mean-cov}
\end{align}
固定选 $k$ 个坐标属于无放回抽样，故
\begin{equation}
  q_2-q_1^2
  =-\frac{k(r-k)}{r^2(r-1)}\le0.
\end{equation}
因此若不作零均值假设，不能只乘一个 $q_2$。

\subsection{适用条件}

真实掩码与梯度都依赖 $Ax$，因此二者并不严格独立。式 \eqref{eq:exact-cov-reduction} 严格对应均匀随机掩码代理模型；真实 top-$k$ 的缩放因子取决于掩码和梯度的联合分布。

\section{跨任务随机子空间近似正交}

\subsection{交叉 Gram 矩阵的均值}

取两个任务独立初始化的随机投影
\begin{equation}
  A_i,A_j\in\R^{r\times n},
  \qquad i\ne j,
\end{equation}
每个元素零均值、方差 $\sigma^2=p/(nr^2)$。定义行空间交叉 Gram 矩阵
\begin{equation}
  M_{ij}=A_iA_j^{\mathsf{T}}\in\R^{r\times r}.
\end{equation}
第 $(u,v)$ 个元素为
\begin{equation}
  (M_{ij})_{uv}
  =\sum_{\ell=1}^{n}(A_i)_{u\ell}(A_j)_{v\ell}.
\end{equation}
两个任务的矩阵独立，故
\begin{align}
  \E[(M_{ij})_{uv}]
  &=\sum_{\ell=1}^{n}
    \E[(A_i)_{u\ell}]
    \E[(A_j)_{v\ell}]\\
  &=0,
\end{align}
从而
\begin{equation}
  \boxed{\E[A_iA_j^{\mathsf{T}}]=0_{r\times r}.}
  \label{eq:mean-orthogonality}
\end{equation}
这是“均值正交”，不是说每次随机初始化后都精确正交。

\subsection{单个元素的方差}

令 $X_\ell=(A_i)_{u\ell}$、$Y_\ell=(A_j)_{v\ell}$。不同任务独立且零均值，交叉坐标项的期望为零，因此
\begin{align}
  \Var((M_{ij})_{uv})
  &=\E\left(\sum_{\ell=1}^{n}X_\ell Y_\ell\right)^2\\
  &=\sum_{\ell=1}^{n}\E[X_\ell^2]\E[Y_\ell^2]\\
  &=n\sigma^4\\
  &=n\left(\frac{p}{nr^2}\right)^2
  =\frac{p^2}{nr^4}.
  \label{eq:cross-entry-var}
\end{align}
于是切比雪夫不等式给出
\begin{equation}
  \Pp\left(|(M_{ij})_{uv}|\ge\epsilon\right)
  \le\frac{p^2}{nr^4\epsilon^2}.
  \label{eq:entry-cheb}
\end{equation}

\subsection{从元素界到矩阵范数界}

若 $\|M_{ij}\|_\F\ge\epsilon r$，则 $r^2$ 个元素中至少有一个的绝对值不小于 $\epsilon$；否则所有元素均小于 $\epsilon$，Frobenius 范数就小于 $\sqrt{r^2\epsilon^2}=\epsilon r$。因此由并合界
\begin{align}
  \Pp(\|M_{ij}\|_\F\ge\epsilon r)
  &\le\sum_{u=1}^{r}\sum_{v=1}^{r}
  \Pp(|(M_{ij})_{uv}|\ge\epsilon)\\
  &\le r^2\frac{p^2}{nr^4\epsilon^2}
  =\frac{p^2}{nr^2\epsilon^2}.
  \label{eq:paper-union-bound}
\end{align}
再由 $\|M\|_2\le\|M\|_\F$，得到论文定理 3.4 的有效但较松界
\begin{equation}
  \boxed{
  \Pp(\|A_iA_j^{\mathsf{T}}\|_2\ge\epsilon r)
  \le\frac{p^2}{nr^2\epsilon^2}.}
  \label{eq:paper-orth-bound}
\end{equation}

还可以直接计算
\begin{align}
  \E\|M_{ij}\|_\F^2
  &=\sum_{u,v}\E[(M_{ij})_{uv}^2]\\
  &=r^2\frac{p^2}{nr^4}
  =\frac{p^2}{nr^2}.
  \label{eq:expected-cross-frob}
\end{align}
对 $\|M_{ij}\|_\F^2$ 使用 Markov 不等式，甚至得到更强的同阈值界
\begin{equation}
  \Pp(\|M_{ij}\|_\F\ge\epsilon r)
  \le\frac{p^2}{nr^4\epsilon^2}.
  \label{eq:stronger-frob-bound}
\end{equation}
式 \eqref{eq:paper-orth-bound} 仍然成立，只是没有利用所有元素平方和的期望。

\subsection{如何传递到 LoRA 更新矩阵}

令两个任务的更新为
\begin{equation}
  U_i=B_iA_i,
  \qquad
  U_j=B_jA_j.
\end{equation}
其 Frobenius 内积逐步变形为
\begin{align}
  \langle U_j,U_i\rangle_\F
  &=\tr(U_j^{\mathsf{T}} U_i)\\
  &=\tr(A_j^{\mathsf{T}} B_j^{\mathsf{T}} B_iA_i)\\
  &=\tr(B_j^{\mathsf{T}} B_iA_iA_j^{\mathsf{T}}).
  \label{eq:update-inner-trace}
\end{align}
最后一步用到了迹的循环不变性：只要乘积维度合法，$\tr(XYZ)=\tr(ZXY)$。

令 $C_{ji}=B_j^{\mathsf{T}} B_i$、$M_{ij}=A_iA_j^{\mathsf{T}}$。Frobenius Cauchy--Schwarz 不等式给出
\begin{equation}
  \boxed{
  |\langle U_j,U_i\rangle_\F|
  =|\tr(C_{ji}M_{ij})|
  \le\|C_{ji}\|_\F\|M_{ij}\|_\F.}
  \label{eq:quantitative-update-orth}
\end{equation}
进一步，
\begin{equation}
  \|B_j^{\mathsf{T}} B_i\|_\F
  \le\|B_j\|_2\|B_i\|_\F,
\end{equation}
所以只有在 $A_iA_j^{\mathsf{T}}$ 小且 $B_i,B_j$ 的范数受到控制时，才能定量地说更新近似正交。原文“对任意学得的 $B_iA_i,B_jA_j$ 都有内积约为零”需要隐含这一范数条件；若任意放大 $B$，绝对内积也可被任意放大。

\subsection{模型合并范数的精确展开}

对 $t$ 个任务，令
\begin{equation}
  \Delta W_{\mathrm{merge}}
  =\sum_{i=1}^{t}w_iU_i,
  \qquad
  W'=W_0+\Delta W_{\mathrm{merge}}.
\end{equation}
由双线性展开，
\begin{align}
  \|\Delta W_{\mathrm{merge}}\|_\F^2
  &=\left\langle\sum_iw_iU_i,\sum_jw_jU_j\right\rangle_\F\\
  &=\sum_{i=1}^{t}w_i^2\|U_i\|_\F^2
   +2\sum_{1\le i<j\le t}w_iw_j\langle U_i,U_j\rangle_\F\\
  &=\sum_{i=1}^{t}w_i^2\|U_i\|_\F^2
   +\sum_{i\ne j}w_iw_j\langle U_i,U_j\rangle_\F.
  \label{eq:merge-exact}
\end{align}
若所有交叉内积相对较小，则
\begin{equation}
  \|\Delta W_{\mathrm{merge}}\|_\F^2
  \approx\sum_{i=1}^{t}w_i^2\|U_i\|_\F^2.
  \label{eq:merge-approx}
\end{equation}
近似误差可以显式控制：
\begin{align}
 &\left|
 \|\Delta W_{\mathrm{merge}}\|_\F^2
 -\sum_iw_i^2\|U_i\|_\F^2
 \right|\\
 &\quad\le
 2\sum_{i<j}|w_iw_j|\,
 \|B_j^{\mathsf{T}} B_i\|_\F\,
 \|A_iA_j^{\mathsf{T}}\|_\F.
 \label{eq:merge-error-bound}
\end{align}
\end{document}
